% Insert after the generic log-canonical component of Paper VIII.
% Requires the v31 notation and the bibliography item LWW.
\section{対数型の共鳴壁と23次元成分}
\label{subsec:resonance23-v32}

一般対数型の非消滅条件を外すと、四次項が五方向同時に自由になる壁が現れる。
その壁は、既知の19次元成分に接する別の23次元成分を作る。
以下ではPoisson構造の53変数のスキームを $\mathscr P$ と書く。

固有座標を $(x,y,z)$ とし、
\[
 p=(\alpha yz,\beta zx,\gamma xy),\qquad
 t=\alpha+\beta+\gamma,\qquad c=txyz
\]
とする。$D=\operatorname{curl}p$ の固有値は
\[
 (\delta_1,\delta_2,\delta_3)
 =(\gamma-\beta,\alpha-\gamma,\beta-\alpha)
\]
である。次の開条件を課す。
\begin{enumerate}
\item $D$ は正則半単純で、$\ker(D|_{S_3})=\mathbb Cxyz$、$t\neq0$。
\item $3(\alpha,\beta,\gamma)-tm\neq0$ がすべての
      $m\in\mathbb Z_{\geq0}^3$、$|m|=3$ に対して成り立つ。
\item $a_2=3\beta+t\neq0$、$a_3=3\gamma+t\neq0$。
\end{enumerate}
この開集合の中で $a=a_1=3\alpha+t$ を考える。
壁 $a=0$ は空でなく、例えば $(\alpha,\beta,\gamma)=(-1,0,4)$ は
$t=3$、$D$ の固有値 $(4,-5,1)$ を与え、上の条件を満たす。

\begin{theorem}[23次元共鳴成分と19次元成分との交叉]
\label{thm:resonance23-v32}
上の開集合において、$a=0$ 上の混合方程式の核はEuler商で15次元である。
三次式による $w$ の移動を行うと、すべての解は一意に
\[
 [F]=[(f(y,z),0,0)],\qquad f\in\mathbb C[y,z]_4
\]
という五次元の残余を持つ標準形になる。
ここで一意性は、混合項を三次式による移動で割った商についての主張である。

壁上の残余 $f=0$ の任意の点で、$\mathscr P$ の完備局所環は
\[
 \mathbb C[[u_1,\ldots,u_{18},a,z_0,\ldots,z_4]]
       /(az_0,\ldots,az_4)
\]
と同型である。したがって、この近傍には既知の19次元対数型成分と
新しい23次元成分の二つだけがあり、両者は滑らかな18次元部分で交わる。
この局所スキームは被約である。

固有座標の置換で得られる三つの共鳴壁の像は、全体として一つの
既約23次元成分 $\mathscr C_{114}$ をなす。
その一般点は $a=0$、$f\neq0$ であり、$\mathscr P$ の滑らかな点である。
\end{theorem}

\begin{proof}
正則半単純線形場の固有値を順序付ける有限\'etale被覆を取り、
固有線束を局所的に自明化する。条件1により
$D(c)=0$ は $c=txyz$ を強制し、その基底は滑らかな9次元である。
これは固有線の配置6次元、跡零固有値2次元、$t$ の1次元からなる。

四次ベクトル場の文字表示 $F_i=x_i x^m v_i$ を用いると、混合方程式は
\[
 \bigl(3(\alpha,\beta,\gamma)-tm\bigr)\cdot v=0
\]
である。非負の10文字ではEuler商は二次元で、条件2によって方程式の階数は1。
従って各文字につき一つの自由変数が残り、計10次元となる。
負の文字は $m_i=-1$ の場合であり、各 $i$ について五つある。
条件3により $i=2,3$ の10変数は零になる。
$i=1$ の五変数はちょうど $F_1\in\mathbb C[y,z]_4$ の係数であり、
その方程式は同じ係数 $a$ を持つ
\[
 az_0=az_1=\cdots=az_4=0
\]
である。異なる文字は異なる六次単項式へ写るため、追加の結合条件はない。

$a$ は9次元基底の滑らかな局所座標の一つである。
残る基底8変数と非負文字の10自由変数を $u_1,\ldots,u_{18}$ と取れば、
表示した完備局所環を得る。イデアルの等式
\[
 (az_0,\ldots,az_4)=(a)\cap(z_0,\ldots,z_4)
\]
により、被約性と二つの枝、交叉の次元がすべて従う。
$a\neq0$ では $z_i=0$ であるから、19次元の枝は既知の一般対数型成分である。
$z_i$ のいずれかが非零なら $a=0$ が強制され、その近傍は滑らかな23次元になる。

非負文字の一自由度は、一般対数型と同じ計算により三次式による $w$ の移動で
すべて消去できる。ここでは $t\neq0$ と条件2を使う。
移動の像には負の文字がないため、残余はちょうど二変数四次式の五次元である。

最後に大域的な既約性を示す。壁は
\[
 (\alpha,\beta,\gamma)=(\alpha,\beta,-4\alpha-\beta)
\]
という二次元平面の非空な開集合である。その上の15次元の混合核はベクトル束である。
これを $\operatorname{GL}_3$ で動かした像は既約であり、一般のファイバーは
固有基底の対角変更に由来する3次元のものである。従って像の次元は
\[
 9+2+15-3=23
\]
となる。固有座標の置換も $\operatorname{GL}_3$ の作用に含まれるので、
三つの壁は同じ既約閉包に入る。
上で得た滑らかな23次元点を含むことから、この閉包は
$\mathscr P$ の一つの既約成分である。
\end{proof}

\begin{remark}[検証と量化]
基準点 $(\alpha,\beta,\gamma)=(-1,0,4)$ では混合写像の階数は20、
Euler商の核は15次元、三次移動の像は10次元である。
$F=0$ における全38方程式のJacobian階数は29、接空間は24次元であり、
$F=(y^4,0,0)$ では階数30、接空間は23次元となる。
これらは有理数上で検証できるが、成分性と局所環の表示の根拠は上の文字ごとの証明である。
被約性の結論はこの共鳴開集合に限られ、全Poisson錐の被約性を意味しない。
\end{remark}

\subsection{重み付き対数型 $L(1,1,4)$ との同定}
\label{subsec:log114-v32}

Poisson構造から得られる基本1形式を
\[
 \eta=\iota_E\iota_P
       (dx\wedge dy\wedge dz\wedge dw)
      =c\,dw-3w(p_1dx+p_2dy+p_3dz)-(x\times F)\cdot dx
\]
とする。Euler場を因子に持つ二重ベクトル場はこの式に寄与しないので、
最後の表示には $w$ に依存しない代表 $P$ を用いてよい。

共鳴標準形 $F=(f(y,z),0,0)$ と $t=-3\alpha$ に対し、
\[
 Q=xw+g(y,z),\qquad
 (\lambda_y,\lambda_z,\lambda_Q)=(-3\beta,-3\gamma,t)
\]
と置くと、
\[
 \eta=yzQ\left(\lambda_y\frac{dy}{y}+
                   \lambda_z\frac{dz}{z}+
                   \lambda_Q\frac{dQ}{Q}\right)
\]
となる必要十分条件は
\[
 f=3\beta g-t\,y\partial_y g
\]
である。重み付きEuler条件は
$\lambda_y+\lambda_z+4\lambda_Q=0$ であり、これは $3\alpha+t=0$ と同じである。
二変数四次単項式 $y^kz^{4-k}$ 上の係数は $3\beta-tk$ なので、
これら五つが非零となる一般点では $g$ が一意に存在する。
例えば $(\alpha,\beta,\gamma)=(-1,-2,6)$ はすべての開条件を満たす。

逆に、一般の次数 $(1,1,4)$ の三因子を持つ重み付き対数型1形式を取る。
一次因子を $y,z$ とし、四次因子の $w$ の係数がこれらと独立である開集合で
その係数を $x$ と選ぶ。四次因子は $xw+H_4(x,y,z)$ と書ける。
$H_4$ の $x$ で割り切れる部分を三次式による $w$ の移動で消すと
$Q=xw+g(y,z)$ となり、上の表示に帰着する。
従って $\mathscr C_{114}$ は、この重み付き対数型 $L(1,1,4)$ の閉包に等しい。
この同定は一般点の幾何的な記述であり、共鳴の追加壁で同じ標準化が一意に
成立するという主張ではない。

